\chapter{Metodologia}
\label{cap:metodologia}

Este capítulo descreve a arquitetura do \textit{framework} proposto, os procedimentos experimentais para sua avaliação e as métricas utilizadas para medir sua eficácia. A metodologia baseia-se em três principais etapas:
\begin{itemize}
    \item Validação empírica das técnicas não invasivas revisadas no Capítulo~\ref{cap:revisao}, com seleção da abordagem que apresentar melhor desempenho segundo \textit{benchmarks} especializados;

    \item Implementação de um \textit{framework} socrático com memória externa persistente, focado em recuperar informações externas ao treinamento do modelo e armazenar o histórico conversacional para busca de contraexemplos;

    \item Avaliação qualitativa e quantitativa da interação dialética produzida.
\end{itemize}

\section{Validação Empírica de Técnicas Não-Invasivas}
\label{sec:validacao-tecnicas}

Para selecionar a melhor estratégia de extensão de \textbf{janela de contexto} para o \textit{framework} proposto, conduzimos uma avaliação de desempenho das técnicas revisadas na Seção~\ref{sec:revisao-tecnicas}, utilizando três \textit{benchmarks} especializados em contexto longo: Needle-in-a-Haystack, RULER e LongBench, conforme justificado também na Seção~\ref{sec:validacao-tecnicas}. Os experimentos foram conduzidos utilizando imagens \textbf{Docker}, código em Python e a API Groq com quatro modelos: GPT-OSS~120B, GPT-OSS~20B, Llama~3.1~8B e Llama~3.3~70B.

\subsection{Estratégias de Avaliação}

As seguintes estratégias foram avaliadas:
\begin{itemize}
    \item \textbf{Raw (Baseline)}: processamento direto do contexto completo até o limite da janela, sem técnica de otimização;

    \item \textbf{Parallel Window}: segmentação em janelas paralelas processadas simultaneamente, seguida de agregação dos resultados, conforme descrito em \cite{ratner2023parallelcontextwindowslarge};

    \item \textbf{Sliding Window}: janela deslizante com sobreposição entre segmentos, conforme descrito em \cite{kamradt2023workaroundtokenlimit};

    \item \textbf{RIG (Dartboard)}: recuperação de informação através de memória externa, priorizando ganho de informação relevante via ranking híbrido semântico-lexical \cite{pickett2024betterragusingrelevantinformationgain};

    \item \textbf{Semantic Compression}: compressão de \textit{prompt} orientada por preservação semântica \cite{gilbert2023semanticcompressionlargelanguage};

    \item \textbf{RSAW} (\textit{Relevance-Stratified Adaptive Window}): estratégia proposta neste trabalho, que combina segmentação baseada em tokens, pontuação Dartboard e estratificação em três níveis de compressão com orçamento de tokens controlado.
\end{itemize}

\subsection{Resultados: Acurácia e Latência}

As Tabelas~\ref{tab:benchmark-results} e~\ref{tab:benchmark-latency} apresentam, respectivamente, as pontuações de acurácia (normalizadas entre 0 e 1) e a latência média (em segundos) obtidas por cada estratégia nos três \textit{benchmarks} selecionados. Os modelos Llama~3.1~8B e Llama~3.3~70B foram executados nos três \textit{benchmarks} completos. Os modelos GPT-OSS foram executados com Needle-in-a-Haystack e RULER em uma primeira rodada e com LongBench em rodada isolada, a fim de evitar esgotamento de cota da API (\textit{rate limit}); as células marcadas com ``---'' indicam ausência de dados nessa configuração.

\begin{table}[htbp]
    \centering
    \caption{Acurácia por estratégia nos \textit{benchmarks} de contexto longo}
    \label{tab:benchmark-results}
    \begin{tabular}{llccc|c}
    \toprule
    \textbf{Modelo} & \textbf{Estratégia} & \textbf{Needle} & \textbf{RULER} & \textbf{LongBench} & \textbf{Média} \\
    \midrule
    \multirow{6}{*}{Llama~3.1~8B}
     & Parallel Window    & 1.00 & 1.00 & 0.95 & 0.98 \\
     & Raw (Baseline)     & 1.00 & 0.94 & 0.95 & 0.96 \\
     & RIG (Dartboard)    & 1.00 & 1.00 & 0.95 & 0.98 \\
     & RSAW (Proposto)    & 0.20 & 0.46 & 0.92 & 0.53 \\
     & Sem. Compression   & 0.33 & 0.56 & 0.75 & 0.55 \\
     & Sliding Window     & 1.00 & 1.00 & 0.95 & 0.98 \\
    \midrule
    \multirow{6}{*}{Llama~3.3~70B}
     & Parallel Window    & 1.00 & 1.00 & 0.95 & 0.98 \\
     & Raw (Baseline)     & 1.00 & 1.00 & 1.00 & 1.00 \\
     & RIG (Dartboard)    & 1.00 & 1.00 & 0.95 & 0.98 \\
     & RSAW (Proposto)    & 0.40 & 0.62 & 0.62 & 0.55 \\
     & Sem. Compression   & 0.33 & 0.69 & 0.70 & 0.57 \\
     & Sliding Window     & 1.00 & 1.00 & 0.95 & 0.98 \\
    \midrule
    \multirow{6}{*}{GPT-OSS~20B}
     & Parallel Window    & 1.00 & 1.00 & 0.92 & 0.97 \\
     & Raw (Baseline)     & 1.00 & 1.00 & 0.83 & 0.94 \\
     & RIG (Dartboard)    & 1.00 & 0.00 & 0.92 & 0.64 \\
     & RSAW (Proposto)    & 0.16 & 0.35 & 0.58 & 0.36 \\
     & Sem. Compression   & 0.80 & 0.00 & 0.33 & 0.38 \\
     & Sliding Window     & 1.00 & 1.00 & 1.00 & 1.00 \\
    \midrule
    \multirow{6}{*}{GPT-OSS~120B}
     & Parallel Window    & 1.00 & 0.72 & 0.67 & 0.80 \\
     & Raw (Baseline)     & 1.00 & 0.64 & 0.62 & 0.75 \\
     & RIG (Dartboard)    & 0.99 & 0.00 & 0.88 & 0.62 \\
     & RSAW (Proposto)    & 0.65 & 0.56 & 0.33 & 0.51 \\
     & Sem. Compression   & 1.00 & 0.13 & 0.62 & 0.58 \\
     & Sliding Window     & 0.87 & 0.81 & 0.79 & 0.82 \\
    \bottomrule
    \end{tabular}
\end{table}

\begin{table}[htbp]
    \centering
    \caption{Latência média por estratégia (segundos)}
    \label{tab:benchmark-latency}
    \begin{tabular}{llccc|c}
    \toprule
    \textbf{Modelo} & \textbf{Estratégia} & \textbf{Needle} & \textbf{RULER} & \textbf{LongBench} & \textbf{Média} \\
    \midrule
    \multirow{6}{*}{Llama~3.1~8B}
     & Parallel Window    & 6.2s & 3.7s &  4.4s &  4.8s \\
     & Raw (Baseline)     & 0.3s & 1.0s &  0.7s &  0.7s \\
     & RIG (Dartboard)    & 6.3s & 3.9s &  6.6s &  5.6s \\
     & RSAW (Proposto)    & 21.2s & 8.4s &  3.2s & 10.9s \\
     & Sem. Compression   & 11.7s & 7.3s & 13.9s & 11.0s \\
     & Sliding Window     & 2.3s & 3.8s &  3.6s &  3.3s \\
    \midrule
    \multirow{6}{*}{Llama~3.3~70B}
     & Parallel Window    & 0.3s & 2.3s &  1.5s &  1.4s \\
     & Raw (Baseline)     & 0.2s & 0.3s &  0.4s &  0.3s \\
     & RIG (Dartboard)    & 3.2s & 2.5s &  2.6s &  2.8s \\
     & RSAW (Proposto)    & 6.0s & 3.2s &  1.3s &  3.5s \\
     & Sem. Compression   & 3.6s & 5.1s &  3.9s &  4.2s \\
     & Sliding Window     & 0.3s & 0.5s &  0.3s &  0.4s \\
    \midrule
    \multirow{6}{*}{GPT-OSS~20B}
     & Parallel Window    & 7.2s & 4.0s &  4.8s &  5.3s \\
     & Raw (Baseline)     & 3.1s & 2.7s &  0.7s &  2.2s \\
     & RIG (Dartboard)    & 7.1s & 0.2s &  4.4s &  3.9s \\
     & RSAW (Proposto)    & 30.6s & 15.5s & 12.1s & 19.4s \\
     & Sem. Compression   & 14.5s & 0.3s & 20.9s & 11.9s \\
     & Sliding Window     & 7.6s & 4.1s &  4.1s &  5.3s \\
    \midrule
    \multirow{6}{*}{GPT-OSS~120B}
     & Parallel Window    & 7.2s & 4.1s &  4.4s &  5.2s \\
     & Raw (Baseline)     & 3.3s & 2.3s &  0.6s &  2.1s \\
     & RIG (Dartboard)    & 7.2s & 0.2s &  3.3s &  3.6s \\
     & RSAW (Proposto)    & 43.2s & 15.8s & 13.7s & 24.2s \\
     & Sem. Compression   & 16.4s & 1.5s & 22.9s & 13.6s \\
     & Sliding Window     & 7.7s & 3.9s &  2.9s &  4.8s \\
    \bottomrule
    \end{tabular}
\end{table}

\subsection{Discussão dos Resultados e Seleção da Estratégia}

Os resultados revelam padrões distintos entre modelos e estratégias. Nos modelos Llama, \textbf{Parallel Window}, \textbf{RIG} e \textbf{Sliding Window} obtiveram acurácia equivalente (0.98), porém \textbf{Sliding Window} destacou-se pela menor latência média (3.3s no Llama~3.1~8B e apenas 0.4s no Llama~3.3~70B), tornando-a a escolha mais eficiente entre as estratégias competitivas. O \textbf{Raw (Baseline)} apresentou desempenho próximo (0.96--1.00), confirmando que, para contextos dentro do limite da janela, o processamento direto é competitivo. \textbf{Semantic Compression} apresentou o pior desempenho geral, com acurácia de apenas 0.55--0.57 nos modelos Llama e degradação severa no RULER para os modelos GPT-OSS.

Nos modelos GPT-OSS, \textbf{Sliding Window} também se destacou, atingindo acurácia de 1.00 tanto no GPT-OSS~20B quanto no GPT-OSS~120B (0.82). O \textbf{RIG (Dartboard)} apresentou colapso no RULER (0.00) para ambos os modelos GPT-OSS, indicando incompatibilidade do mecanismo de recuperação com o formato de perguntas sintéticas desse \textit{benchmark}.

A estratégia \textbf{RSAW} proposta neste trabalho apresentou resultados mistos: embora não supere as estratégias mais simples nos \textit{benchmarks} de recuperação direta (Needle e RULER), demonstrou maior equilíbrio entre os três \textit{benchmarks}, sem zerar em nenhum deles. Sua maior limitação é a latência elevada (10.9s--24.2s de média), decorrente do pipeline de compressão em múltiplas etapas.

Com base nesses resultados, \textbf{Sliding Window} foi selecionada como estratégia base para o \textit{framework} proposto, dada a combinação de alta acurácia, baixa latência e comportamento estável entre modelos e \textit{benchmarks}. A estratégia RSAW permanece como contribuição experimental, com potencial para cenários onde o equilíbrio entre \textit{benchmarks} heterogêneos é prioritário sobre latência.
